% ============================================================
\section{Introduction}
\label{sec:intro}
% ============================================================

Network-on-Chip (NoC) has become the dominant communication paradigm for multi-core processors, AI accelerators~\cite{journeys}, and System-on-Chips (SoCs). As the number of integrated cores scales from tens to hundreds~\cite{amd2023,apple2023}, the efficiency of the on-chip interconnection network directly determines overall system performance in terms of latency, throughput, and energy consumption~\cite{dally2004}.

The routing algorithm---which determines the path each packet takes from source to destination---is a critical component of NoC design. Conventional deterministic algorithms such as XY dimension-order routing (DOR)~\cite{dally2004} are widely used due to their simplicity and deadlock freedom. However, deterministic routing always selects the same path for each source-destination pair regardless of network congestion, leading to structural bottlenecks.

Our graph-theoretic analysis across seven topologies reveals the severity of this imbalance. Under XY routing on an $8\times8$ mesh, the congestion imbalance---measured as the coefficient of variation ($\sigma/\mu$) of link utilization---reaches 0.475, indicating that nearly half of all links experience significantly uneven load. Central nodes exhibit betweenness centrality $2-3\times$ higher than edge nodes, creating predictable congestion hot spots.

Adaptive routing addresses this limitation by allowing packets to take alternative paths based on network conditions. Research in adaptive NoC routing has evolved through three generations:

\begin{enumerate}
    \item \textbf{Heuristic adaptive routing}~\cite{dyad2004,ebrahimi2019} uses local congestion thresholds to dynamically select output ports. While simple to implement, these approaches rely on hand-crafted heuristics that cannot capture global traffic patterns.
    
    \item \textbf{Deep Reinforcement Learning (DRL) routing}~\cite{wang2022maar,deepnr2022,garnn2023} employs neural network agents trained via RL to make routing decisions. These methods achieve 20--35\% latency improvement over deterministic routing. However, they use Multi-Layer Perceptron (MLP) encoders that process flattened router states as vectors, discarding the inherent graph structure of the NoC.
    
    \item \textbf{Graph Neural Network (GNN) routing for wide-area networks}~\cite{g-routing2023,drl-gnn2019} applies message-passing neural networks to network routing problems. GNNs naturally operate on graph-structured data and can capture topological relationships. However, these approaches target software-defined networks with millisecond-scale latency requirements, not the nanosecond-scale constraints of on-chip routing.
\end{enumerate}

The fundamental challenge in applying GNN-based routing to NoC is \textbf{inference latency}. A single GNN forward pass on a 64-node graph requires approximately $10^6$ cycles on a CPU, while each NoC routing decision must complete within $5-10$ cycles~\cite{wang2022maar}. This approximately five-order-of-magnitude gap renders per-packet GNN inference infeasible.

\textbf{Our approach:} GNNocRoute-DRL addresses this challenge through \textbf{periodic policy optimization}. Instead of making per-packet routing decisions with GNNs, we propose a two-timescale architecture: (1) a \textit{slow timescale} where GNN re-encoding runs sporadically (every $10^5$--$10^6$ cycles) in the background to refresh topology-aware node embeddings, while lightweight DRL policy updates occur every $5,000$ cycles; and (2) a \textit{fast timescale} where per-packet routing decisions use pre-computed Q-values and local congestion information via table lookup.

\subsection{Contributions}

This paper makes the following contributions:

\begin{enumerate}
    \item \textbf{Congestion analysis:} Quantitative demonstration across seven topologies that deterministic routing creates structurally predictable bottlenecks, with congestion imbalance of 0.475 on $8\times8$ mesh.
    
    \item \textbf{GNN-topology correlation:} Experimental evidence showing GATv2 node embeddings correlate with betweenness centrality ($|r| = 0.978$), validating that topology-aware representations emerge implicitly.
    
    \item \textbf{GNNocRoute-DRL framework:} A complete architecture combining GNN-based topology encoding with DRL-based routing decisions, designed for periodic operation under NoC timing constraints.
    
    \item \textbf{Comprehensive evaluation:} Cycle-accurate simulation (BookSim2, 252 configurations) comparing four routing algorithms across three topologies and three traffic patterns, demonstrating up to 16.9\% latency reduction for hotspot traffic on $8\times8$ mesh.
    
    \item \textbf{Cross-topology transfer:} Demonstration that GNN embeddings pre-trained on small topologies transfer effectively to larger configurations (Pearson $r \approx 0.96$), reducing re-training cost.
\end{enumerate}
